% ===== §25 Contributions and Structure =====
\section{Contributions and the Structure of the Paper}
\label{p22:sec:contributions}

\noindent
This closing section states the paper's contributions explicitly and exhibits its structure in a form that the running argument, of necessity linear, could not make visible all at once. The visualizations are not decorations but the connective tissue the prose approached one link at a time; gathered here, they show the joints between claims that the linear exposition could only assert in sequence. A word on the spirit in which they are offered is owed first, because it governs how the tables that compare theories are to be read. The paper holds to a dialectical-materialist commitment, which means that it treats no theory gathered here, neither the predictive nor the contemplative nor the phenomenological nor the psychoanalytic, as a final or unconditioned rationality, and assigns each instead its determinate place, its proper level, and its limit. The comparative tables therefore record, for each theory, not only what it contributes but what it cannot reach, and the layered syntheses they summarize are arrangements in which each theory does the work it can and is kept from the work its commitments would botch. This is the materialist refusal of any theory's pretension to be the whole, carried into the form of the exposition itself.

\subsection{The reduction chain}
\label{p22:sub:contrib-chain}

The argument's spine is a reduction chain, and the chain is best seen whole. Each link carries a determinate argumentative burden, discharged in a determinate section, and the chain's force is cumulative, each link presupposing the one before and enabling the one after.

\begin{figure}[ht]
\centering
\resizebox{\textwidth}{!}{%
\begin{tikzpicture}[node distance=0.55cm and 0.85cm]
\node[chainbox] (rp) {Relational production\\ \footnotesize(\S\ref{p22:sec:field-ontology})};
\node[chainbox, right=of rp] (field) {requires a \textbf{field}\\ \footnotesize a conditional environment of coupling};
\node[chainbox, right=of field] (aes) {whose construction is \textbf{aesthetic}\\ \footnotesize not engineering (\S\ref{p22:sec:not-engineering})};
\node[chainbox, below=of aes] (eth) {the aesthetic is coupled to the \textbf{ethical}\\ \footnotesize one perception (\S\ref{p22:sec:coupling})};
\node[chainbox, left=of eth] (edu) {so field-building is \textbf{aesthetic education}\\ \footnotesize the terminus (\S\ref{p22:sec:terminus})};
\node[chainbox, left=of edu] (cult) {the \textbf{cultivation}\\ \footnotesize perception, practice, maintenance; the four-step model};
\draw[chainarrow] (rp) -- (field);
\draw[chainarrow] (field) -- (aes);
\draw[chainarrow] (aes) -- (eth);
\draw[chainarrow] (eth) -- (edu);
\draw[chainarrow] (edu) -- (cult);
\end{tikzpicture}}
\caption{The reduction chain. Each link is argued in the section noted; the orthogonality guard (\S\ref{p22:sec:orthogonality}) bounds the third link, marking where the aesthetic-ethical coupling ends and the separate question of justice begins.}
\label{p22:fig:chain}
\end{figure}

The chain is not a deduction in which the conclusion is contained in the premises, but a reduction in which each link is independently argued and the terminus reframes what aesthetic education is. The most contested link is the second, the claim that field-building is aesthetic and not an engineering matter, and it bears the weight of the whole, since if the construction of a field could be specified by rule the chain would terminate in a technics rather than an education. The orthogonality guard is the chain's necessary qualification: the coupling of the aesthetic and the ethical holds at the level of the field's generative direction and is silent on the justice of who is rendered fuel for the accumulation, and this silence, marked at the third link, is what keeps the chain from issuing in an aestheticism.

\subsection{The four-step model, mapped}
\label{p22:sub:contrib-fourstep}

The four-step model is the cultivation in motion, and its four movements stand in determinate relations to the relational state that makes each the principal site, to the failure that shadows each, and to the Eastern resource that most illuminates each. Table~\ref{p22:tab:fourstep} gathers these relations, which the prose introduced one movement at a time.

\begin{table}[ht]
\centering
\small
\begin{tabularx}{\textwidth}{@{}L{2.1cm} L{2.9cm} L{3.0cm} L{3.0cm}@{}}
\toprule
\thd{Movement} & \thd{Principal when the field is\dots} & \thd{Shadowed by the failure of\dots} & \thd{Eastern resource} \\
\midrule
Perception & dull (sight has gone slack) & the flattened morning (prior absorbs the contingency) & \makecell[tl]{single-practice\\absorption;\\emptying-to-await} \\
Reception & lately wounded (defence is high) & defensive gating; or grasping possession & the mirror that responds yet does not store \\
Sharing & estranged (beauty does not cross) & hoarding; or over-speaking that crushes & music as the common; concord without uniformity \\
Reproduction & stalled in plenty (consumed not thickened) & consumption demanding ever-larger novelty & daily renewal inscribed on the washing-basin \\
\bottomrule
\end{tabularx}
\caption{The four-step model under the generative dialectic. The principal site migrates with the field's state; the cultivation has no fixed prescription but reads which movement is principal, which reading is itself the first movement turned on the relation as a whole.}
\label{p22:tab:fourstep}
\end{table}

The table makes visible what the prose asserted but could not display, that the four movements are not a procedure performed in fixed order but a diagnosis: which movement is principal is read from the field's state, and the right labour is a function of that reading rather than a standing rule. The migration of the principal site across the four relational states is the practical content of the generative dialectic, and it is what distinguishes this account from the issuing of four techniques to be applied uniformly.

\subsection{The four hardest steps, unified}
\label{p22:sub:contrib-hardest}

At each of the four theoretical battlefields, a conflict threatened the model and was resolved by an inversion of the same shape, the shape in which a relinquishing of grasp founds rather than forbids a generation. Table~\ref{p22:tab:hardest} sets the four inversions side by side, exhibiting the single movement of thought that the prose traced four times in four idioms.

\begin{table}[ht]
\centering
\small
\begin{tabularx}{\textwidth}{@{}L{2.0cm} L{3.4cm} L{3.4cm} L{2.3cm}@{}}
\toprule
\thd{Movement} & \thd{What is relinquished} & \thd{What it thereby grounds} & \thd{Section} \\
\midrule
Perception & deluded, subject-against-object division (and grasping investment) & the subtle discernment of appropriateness (and open caring) & \S\ref{p22:sec:perception-field} \\
Reception & possessive retention; responding yet not storing & healthy reproduction (curvature, not hoarded stock) & \S\ref{p22:sec:reception} \\
Sharing & the will to fusion, the dissolution of the two & union as concord, which preserves the difference it needs & \S\ref{p22:sec:sharing} \\
Reproduction & the solitary feat of accumulation & the cycle, fed and renewed by the living prior three steps & \S\ref{p22:sec:reproduction} \\
\bottomrule
\end{tabularx}
\caption{The four hardest steps as one movement of thought. Each conflict inverts on the same distinction, between a grasp that possesses and a generation that does not, which is the ancient counsel to generate without possessing read into perception, welcome, the between, and time.}
\label{p22:tab:hardest}
\end{table}

That the four inversions share a structure is the paper's strongest evidence that the four movements are not an arbitrary list but the articulation of one thing. The unity is not imposed by the table; it was earned at each battlefield separately, and the table only makes visible that what was earned four times was, each time, the same.

\subsection{Three beauties and four movements}
\label{p22:sub:contrib-map}

The capacity admits two true and non-competing articulations, the static anatomy of three beauties and the dynamic cycle of four movements. Their relation is a non-trivial mapping, three onto four, which Figure~\ref{p22:fig:map} displays.

\begin{figure}[ht]
\centering
\begin{tikzpicture}[node distance=0.5cm and 2.2cm]
\node[facenode] (perc) {Perception\\ \footnotesize(the input face)};
\node[facenode, below=0.7cm of perc] (prac) {Practice\\ \footnotesize(the active face)};
\node[facenode, below=0.7cm of prac] (main) {Maintenance\\ \footnotesize(the temporal face)};
\node[stepnode, right=of perc] (s1) {Perception};
\node[stepnode, below=0.5cm of s1] (s2) {Reception};
\node[stepnode, below=0.5cm of s2] (s3) {Sharing};
\node[stepnode, below=0.5cm of s3] (s4) {Reproduction};
\node[above=0.15cm of perc, font=\footnotesize\itshape\color{deeprose}] {Three beauties (anatomy)};
\node[above=0.62cm of s1, font=\footnotesize\itshape\color{rose}] {Four movements (cycle)};
\draw[maplink] (perc) -- (s1);
\draw[maplink] (prac) -- (s2);
\draw[maplink] (prac) -- (s3);
\draw[maplink] (main) -- (s4);
\end{tikzpicture}
\caption{The static anatomy and the dynamic cycle of one capacity. The middle face, practice, divides in motion into reception and sharing, the doing that takes the contingency in and carries it across; this is why the anatomy has three faces and the cycle four movements. Neither articulation subsumes the other; the paper holds both, as it holds the plural without forcing it into a governed unity.}
\label{p22:fig:map}
\end{figure}

\subsection{The genealogy, gathered}
\label{p22:sub:contrib-genealogy}

The coupling of the aesthetic and the ethical is an inheritance recognized across traditions, not the paper's stipulation, and Table~\ref{p22:tab:genealogy} gathers the witnesses with the proposition each contributes and the link of the paper's argument each supports. The table also records the asymmetry between the Western witnesses, for whom the coupling is a recovery against their own separation, and the Eastern witnesses, for whom it is an inheritance never lost.

\begin{table}[ht]
\centering
\small
\begin{tabularx}{\textwidth}{@{}L{2.3cm} L{4.6cm} L{2.0cm} L{1.5cm}@{}}
\toprule
\thd{Witness} & \thd{The recognition contributed} & \thd{Supports} & \thd{Mode} \\
\midrule
Rite-and-music (Confucian) & the person is completed in the aesthetic; the good and beautiful never divided & the coupling itself & inheritance \\
Schiller & the moral problem must first be solved as an aesthetic one & aesthetic education grounds the ethical & recovery \\
Wittgenstein & ethics and aesthetics are one; both shown, not said & the appropriate overflows rules & recovery \\
Scheler & value is felt before it is reasoned & value-apprehension is perceptual (\S\ref{p22:sec:coupling}) & recovery \\
Gadamer & taste was a moral power before it was narrowed to art & the fitting where no rule applies & recovery \\
Williams & moral theory mistakes the nature of ethical life & appropriateness is not rule-governed & recovery \\
Murdoch & just attention is moral action; the moral and aesthetic gaze are one & perception is at once aesthetic and moral & recovery \\
\bottomrule
\end{tabularx}
\caption{The genealogy of the aesthetic-ethical coupling. The witnesses share almost nothing in method or metaphysics, which is why their convergence on one recognition is strong evidence; the paper supplies the mechanism, holonomy and the field, that the tradition reported but did not explain.}
\label{p22:tab:genealogy}
\end{table}

\subsection{Contributions}
\label{p22:sub:contrib-list}

The paper's contributions, stated plainly, are these. First, it establishes the concept of \emph{relational aesthetics}, on which beauty is a coupling event between subjects, produced relationally, and it shows this concept to resolve the old quarrel between objectivist and subjectivist aesthetics by locating beauty where both were reaching, in the between. Second, it proposes the \emph{four-step model} of perception, reception, sharing, and reproduction, and shows it to be not a procedure but the smallest form a good cycle takes between two people, a holonomy cycle returning to its root and not its origin. Third, it reconstructs the dialectic of principal and secondary contradiction into a \emph{generative dialectic} that directs generative attention rather than concentrating force, consistent with a dialectical-materialist refusal to absolutize any theory. Fourth, it supplies the \emph{mechanism} of the aesthetic-ethical coupling that the philosophical tradition reported but never explained, the identification of appropriateness with the perceptible signature of accumulating holonomy. Fifth, it conducts four \emph{adversarial syntheses} of the contending traditions of perception, reception, sharing, and reproduction, assigning each theory its level, its reach, and its limit, and resolving their conflicts by inversions that share one structure. Sixth, it installs the \emph{orthogonality guard} that separates the refinement of the sensibility from the justice of the field, securing the account against aestheticism. Seventh, it specifies an \emph{implementation-level interface} of observable signatures, including inter-brain coupling as a candidate correlate of the central concept, without reducing the relational to the neural. Eighth, it gives a \emph{political economy of the field} that defends the cultivation against the charge of leisure-dependence and shows it most needed by the dispossessed. And ninth, it \emph{unifies the series}, naming the cultivable capacity that the prior studies of poetry, the vow, water, and travel each presupposed without naming.

\subsection{The structure of the paper}
\label{p22:sub:contrib-structure}

The paper falls into five parts. Part~I (\S\ref{p22:sec:introduction}--\S\ref{p22:sec:alexander}) lays the foundations: the introduction, the generative dialectic that is the method, the ontology of the field, the argument that field-building is not engineering, and the design-theoretic support. Part~II (\S\ref{p22:sec:genealogy}--\S\ref{p22:sec:orthogonality}) establishes the legitimacy of the aesthetic-ethical coupling: the genealogy, the coupling mechanism, and the orthogonality guard. Part~III (\S\ref{p22:sec:three-beauties}--\S\ref{p22:sec:maintenance}) sets out the three beauties, the static anatomy of the capacity, with perception receiving the densest treatment as the first battlefield. Part~IV (\S\ref{p22:sec:relational-aesthetics}--\S\ref{p22:sec:worked-example}) develops the four-step model and its theory: the establishment of relational aesthetics, the four battlefields of reception, sharing, and reproduction joined to the perception battlefield of Part~III, the implementation-level interface, and the worked example of contingency. Part~V (\S\ref{p22:sec:fetishism}--\S\ref{p22:sec:conclusion}) draws the consequences: the critique of aestheticized consumption, the account of micro self-legislation against the symbolic crisis, the political economy of the field, the terminus that unifies the series, the account of transmission, and the polyphonic conclusion. The present section stands outside the five parts as their retrospective gathering, and the conclusion that follows returns the whole to the everyday from which it began.
